\documentclass{article}
\usepackage{amsmath,amssymb}
\usepackage{booktabs}
\usepackage[margin=1in]{geometry}
\title{Step 231: Ramanujan Tau Carrier and Dichotomy Classification}
\author{Six Birds Foundations III}
\date{2026-05-16}
\begin{document}
\maketitle

\section{Carrier}

Ramanujan's tau function is defined by the modular discriminant
\[
  \Delta(z)=q\prod_{n\ge 1}(1-q^n)^{24}
  =\sum_{n\ge 1}\tau(n)q^n,
  \qquad q=e^{2\pi iz}.
\]
It is the normalized cusp form of weight \(12\) and level \(1\).

The associated Dirichlet series is
\[
  L(\tau,s)=\sum_{n\ge 1}\tau(n)n^{-s}.
\]
The completed form may be written
\[
  \Lambda_\tau(s)=(2\pi)^{-s}\Gamma(s)L(\tau,s),
\]
with functional equation pairing \(s\) and \(12-s\). Thus the critical line is
\[
  \Re(s)=6.
\]

\section{Ramanujan--Petersson Versus RH}

The Ramanujan--Petersson assertion for \(\tau\) controls the size of Fourier coefficients, for example
\[
  |\tau(p)|\le 2p^{11/2}.
\]
Deligne proved this coefficient bound using the Weil conjectures. This theorem is not the same as the RH analog for \(L(\tau,s)\).

The \(L(\tau,s)\)-RH analog is:
\[
  \text{all nontrivial zeros of }L(\tau,s)\text{ lie on }\Re(s)=6.
\]
This is open.

\section{Dichotomy Audit}

The carrier is not CRE for Riemann RH: the zero-line conjecture for \(L(\tau,s)\) is a separate automorphic \(L\)-function RH analog. It is also not natively closed. Hence it is outside the current RH Dichotomy, which classifies CRE carriers and non-CRE carriers with proved native RH analogs.

\section{Generalization Note}

If the framework is generalized to a Selberg-class or automorphic-\(L\)-function RH analog family, the tau carrier would be target-equivalent to its own conjecture. In that generalized setting it would behave like a `CRCFT-TE' instance relative to the target \(L(\tau,s)\)-RH.

\section{Verdict}

\[
  \boxed{\texttt{V\_tau\_outside\_dichotomy}.}
\]

\end{document}
